\section{Introduction}
\label{sec:intro}

Modern deep learning success assumes abundant, diverse training data, yet many high-value applications -- rare disease diagnosis, endangered-species monitoring, low-resource robotics -- can only offer a handful of labeled examples per class \citep{fawakherji2026learning,wang2020generalizing}. Few-shot learning addresses this bottleneck, and among its solutions, gradient-based meta-learning \citep{schmidhuber1987evolutionary,bengio2013optimization}, epitomized by MAML \citep{finn2017model}, dominates precisely because it is architecture-agnostic: a shared initialization $\theta_0$ is meta-trained so a handful of gradient steps on a new task's \emph{support set} $S_i$ suffice to adapt it.

That adaptation, however, is a black box. What prior knowledge in $\theta_0$ drives the few gradient steps taken on $S_i$, and whether that drive helps or hurts, is not observable from the outside. This is not academic: when $S_i$ is hard, noisy, or drawn from a distribution the meta-training tasks never anticipated, adaptation can actively \emph{degrade} query performance -- \emph{negative transfer} \citep{yue2020interventional,si2025meta}. Existing accounts of why MAML adapts are themselves in tension: \citet{raghu2019rapid} argue most of its benefit comes from reusing near-frozen backbone features and only re-fitting the head (ANIL), whereas \citet{oh2020boil} show that under domain shift the opposite regime, backbone-only adaptation, dominates (BOIL). Both accounts, and the negative-transfer literature built on them \citep{yao2019hierarchically,yao2020automated,chen2021hetmaml,parast2026hsfm}, characterize adaptation only through its \emph{symptom} -- a change in accuracy or in an intervened representation -- not, at the level of individual support-image regions, \emph{where in feature space} it went right or wrong.

Meanwhile, post-hoc XAI has matured for static, converged models: Grad-CAM \citep{selvaraju2020grad}, SHAP \citep{lundberg2017unified}, and related attribution methods \citep{simonyan2014visualising,sundararajan2017axiomatic,ribeiro2016lime} all extract a saliency signal from \emph{one} fixed parameter state -- precisely what a bi-level optimizer like MAML does not have, since its meta-test behavior is the endpoint of a $T$-step trajectory $\theta_0 \!\to\! \phi_1 \!\to\! \cdots \!\to\! \phi_T$ that none of these methods aggregate across. The one method reasoning about meta-training influence at the task level, TL-XML \citep{mitsuka2025tlxml}, requires the full meta-training corpus and still treats the adaptation step as an opaque unit -- infeasible when that corpus is proprietary or discarded, and uninformative about \emph{which pixels} of $S_i$ mattered.

This gap motivates two questions: \emph{(i)} for a given target task, is the adaptation step itself helpful or harmful, and by how much; and \emph{(ii)} which support-set features, in the shared feature space $F$ that $\theta_0$ and the adapted parameters both live in, are responsible? We answer \emph{(i)} with \textbf{adaptation gain} $\Delta M$, a differentiable, sign-aware scalar contrasting full adaptation against a feature-reuse baseline that freezes the backbone; and \emph{(ii)} with \textbf{\fama} (Feature-space Adjoint Meta-learning Attribution), which back-propagates $\Delta M$ through the \emph{entire} trajectory via an adjoint recursion, avoiding the $\mathcal{O}(TP^2)$ cost of literal BPTT. Since $S_i$ is available at meta-\emph{test} time by construction, \fama needs no access to the meta-training set.

In summary, our contributions are:
\begin{itemize}[leftmargin=*,itemsep=1pt,topsep=2pt]
    \item \textbf{Adaptation gain $\Delta M$} (\Cref{sec:adaptation_gain}), a differentiable measure of the benefit or harm caused by letting the backbone adapt on $S_i$, which we prove reduces, to leading order in the fast-adaptation regime, to a term that depends \emph{only} on the backbone's displacement (\Cref{prop:decomposition}).
    \item \textbf{\fama} (\Cref{sec:fama}), a post-hoc explainer that extends Grad-CAM across a multi-step adaptation trajectory via a Pearlmutter Hessian-vector-product adjoint recursion, at $\mathcal{O}(TP)$ cost, requiring no meta-training data and no architectural change.
    \item Two \textbf{trajectory-aware faithfulness protocols} (\Cref{sec:protocols}) -- a bidirectional deletion test and cascading-randomization / support-intervention sanity checks -- adapting static-model evaluation principles \citep{samek2016evaluating,petsiuk2018rise,adebayo2018sanity} to a setting where the object of explanation is a trajectory; neither is specific to \fama, so both apply to any future explainer for fast adaptation.
    \item An empirical study on MAML/Conv4 (\Cref{sec:experiments}) that validates \fama quantitatively on tieredImageNet, and, on 600 miniImageNet/tieredImageNet$\to$CUB-200-2011 cross-domain tasks, surfaces \emph{contextual bias} as a concrete, visualizable candidate driver of negative transfer.
\end{itemize}
